% !TeX program = pdflatex
% !TeX encoding = UTF-8
% AsteroidAnchor: the single confirmed, problem-driven mechanism paper.
% Planning manuscript only. No failure mechanism, new design, simulation result,
% experimental benefit, or asteroid-environment performance is established yet.
%
% Close the chain: documented difficulty -> conditions and cause -> targeted
% structural/action change -> model and parameter choice -> matched comparison
% and ablation -> validity range, costs, and failure boundary.
% Theory, evaluation, state recognition, and adaptive actions are supporting
% methods only when the selected problem requires them.
%
% Before drafting claims, freeze this sentence:
% [Reference mechanism] under [terrain/contact/loading conditions] fails to
% meet [mission requirement] because of [testable cause]. Evidence: [source].
% Still uncertain: [alternative explanations and untested conditions].

\documentclass[11pt]{article}
\usepackage[T1]{fontenc}
\usepackage[utf8]{inputenc}
\usepackage{amsmath,amssymb}
\usepackage{booktabs}
\usepackage{graphicx}
\usepackage[hidelinks]{hyperref}
\usepackage[margin=1in]{geometry}
\setlength{\emergencystretch}{2em}

\title{A Problem-Driven Anchoring Mechanism for Rubble-Pile Asteroids:\\
Design, Mechanics, and Validation}
\author{Authors to be confirmed}
\date{Planning manuscript --- problem and evidence pending}

\newcommand{\planned}[1]{\par\noindent\textit{Planned work: #1}\par}
\newcommand{\figureplaceholder}[1]{%
  \fbox{\parbox{0.91\linewidth}{\centering
  \vspace{0.018\textheight}\textbf{Planned evidence --- no results yet}\par
  \smallskip #1\par\vspace{0.018\textheight}}}}
\newif\ifBibliographyReady
\BibliographyReadyfalse

\begin{document}
\maketitle

\begin{abstract}
This is a planning manuscript for one mechanism-centered study of anchoring on
rubble-pile asteroid surfaces. A specific unresolved difficulty, reference
mechanism, and task condition will first be selected from verified literature
and reproduction. The resulting design will target the diagnosed cause and be
tested against the reference under matched mission requirements, including
installation, post-installation load capacity, resource use, and failure modes.
The mechanism and its performance have not yet been determined.
\end{abstract}

\section{Introduction: a documented anchoring problem}
% Separate observed failures, untested conditions, authors' suggestions,
% improvements already implemented, and subsequent publications.
% Keep only literature needed to establish the selected problem and novelty.
\planned{Review original and follow-up studies, choose one task-relevant and
testable difficulty, and state the reference mechanism, terrain/contact/loading
conditions, mission requirement, observed evidence, and remaining uncertainty.}
\planned{Screen the current primary terrain hypothesis of few-pascal, weakly
confined regolith without a fixed support. Treat 0.1, 1, 10, and 100 N as
exploratory task-force levels, not frozen requirements. Define the strength
quantity, moments, duration, allowable motion, and resource limits before
selecting the comparison task or claiming feasibility.}

\begin{figure}[htbp]
\centering
\figureplaceholder{M-F01: Evidence map linking the selected failure, its test
conditions, original figures/data, later fixes, and the unresolved gap.}
\caption{Planned evidence and problem definition; no gap is established yet.}
\end{figure}

\section{Reference mechanism and cause of the difficulty}
% Reproduce one decisive result where feasible. Preserve reference geometry,
% boundary support, material, contact, loading, and reported uncertainty.
% Compare observed failure modes with competing causal explanations.
\planned{Reproduce or independently check a key reference result. Identify when
the difficulty appears, examine the load and reaction paths, and test the most
plausible cause against alternatives before proposing a new design.}

\begin{figure}[htbp]
\centering
\figureplaceholder{M-F02: Reference mechanism, installation and load paths,
reproduced key result, onset conditions, and alternative cause tests.}
\caption{Planned reference and causal diagnosis.}
\end{figure}

\section{Principle and design of the new mechanism}
% Choose a change to structure or action that specifically interrupts the
% diagnosed failure. Compare with the closest prior art at feature level.
% Staging, bio-inspired geometry, sensing, and adaptive deployment are options.
% Show where reaction force and torque close.
\planned{Screen a small number of targeted concepts with minimal tests. Describe
the chosen geometry, actuation sequence, reaction path, mass and power demands,
and exact difference from the closest existing mechanism. State why the change
should affect the measured failure.}

\begin{figure}[htbp]
\centering
\figureplaceholder{M-F03: New geometry, operation sequence, free-body diagrams,
and the single structural or action change being tested.}
\caption{Planned mechanism and its causal hypothesis.}
\end{figure}

\section{Mechanical model and parameter design}
% Use only theory needed to explain this mechanism and choose its parameters.
% Verify contact laws, numerical convergence and boundary conditions; calibrate
% against independent component/material data, not final validation cases.
\planned{Build and check the simplest model able to predict the selected
difficulty and design change. Calibrate critical contact/material properties,
identify parameter sensitivity, and select design values before final tests.}

\begin{figure}[htbp]
\centering
\figureplaceholder{M-F04: Model checks, calibrated quantities, predicted response,
parameter sensitivities, and design choices with uncertainty.}
\caption{Planned model credibility and parameter rationale.}
\end{figure}

\section{Matched comparison and mechanism tests}
% Select the primary metric from the selected problem: installation reaction,
% post-installation holding/deflection, deployment-plus-proof-load success, or
% sensitivity to terrain. In every case test both installation and holding.
% Match mission load, terrain, initial contact, support, allowed attempts, and
% resource budgets. Count every attempt. Compare reference and ablated design.
% Independent units are complete runs or separately prepared beds.
\planned{Freeze task conditions, main metric, minimum meaningful improvement,
failure classes, resource limits, repeats, and exclusions. Compare the reference,
new mechanism, and a suitable ablation using independent simulation and physical
tests. Report all attempts, uncertainty, costs, and failed outcomes.}

\begin{figure}[htbp]
\centering
\figureplaceholder{M-F05: Matched reference--new--ablation comparison of the
selected primary metric, installation, holding, resources, and failure modes.}
\caption{Planned evidence that the proposed change causes a useful improvement.}
\end{figure}

\section{Discussion, limits, and conclusion}
% A better result in one condition does not establish universal adaptability.
% Examine terrain, initial pose, load direction, duration, and low-gravity
% transfer only to the extent supported by model and feasible tests.
\planned{Map the conditions where the new mechanism helps, fails, or remains
untested. Explain disagreements between prediction and experiment. Bound claims
about vacuum and reduced gravity by the actual validation conditions.}

\begin{figure}[htbp]
\centering
\figureplaceholder{M-F06: Applicable and failed regimes, performance--resource
tradeoffs, and remaining extrapolation to asteroid conditions.}
\caption{Planned validity and failure boundaries.}
\end{figure}

\section*{Completion criterion}
\planned{Submit as one mechanism paper only when the problem is documented,
the cause is supported, the design targets that cause, and matched independent
comparisons establish the effect and its limits. Simulation setup, generic
parameter scans, evaluation scores, and test-rig construction are supporting
work unless they later produce a separately validated research question.}

\section*{Data and software availability}
\planned{Archive the screened literature evidence, reproduction inputs,
model/calibration settings, CAD and prototype revisions, complete-run raw data,
failures, and scripts that produce the final figures.}

\ifBibliographyReady
  \bibliographystyle{plain}
  \bibliography{references}
\else
  \section*{References to verify}
  Add only original publications and follow-up studies checked against their
  source text. Bibliography output is disabled at the planning stage.
\fi

\end{document}
